\chapter{Inversion}\label{ch:inversion}

\begin{center}\itshape
Four circles to the kissing come.\\
The smaller are the benter.\\
The bend is just the inverse of\\
The distance from the center.
\upshape(Frederick Soddy, \emph{The Kiss Precise}, Nature, 1936)
\end{center}

In November 1643 Ren\'e Descartes wrote to Princess Elisabeth of Bohemia, his most gifted correspondent, about an old problem of
Apollonius: given three circles, find a circle tangent to all three. Elisabeth had attacked it with algebra and been buried in
square roots. Descartes answered with a relation between the radii of four mutually tangent circles, but he wrote it in a form so
complicated that nobody noticed how simple it was. Almost three centuries later the chemist Frederick Soddy, who had won the 1921 Nobel
Prize for his work on isotopes, rediscovered the relation and published it in \emph{Nature} in 1936 as a poem, \emph{The Kiss
Precise}, whose second stanza opens with the lines of the epigraph. In modern form it is one line: if four circles touch each other, their
curvatures $k_i=1/r_i$ satisfy
\[
(k_1+k_2+k_3+k_4)^2=2\,(k_1^2+k_2^2+k_3^2+k_4^2).
\]

This chapter explains Descartes' theorem, and several other theorems about tangent circles, with one transformation: inversion in a
circle. It was used by Jakob Steiner in the 1820s and systematised by Pl\"ucker, Liouville and William Thomson (Lord Kelvin), who used
it in 1845 to solve problems in electrostatics. Inversion turns circles into lines and lines into circles, keeps angles, and so it
can turn a hard configuration of tangent circles into an easy one of parallel lines or concentric circles, solve it there, and bring the
answer back. It is also the key to the M\"obius transformations of the complex plane, which are the symmetries of the hyperbolic plane of
Chapter~10.

\begin{think}
Draw a circle inside another one, not touching it. Fit a circle between them, touching both; then a second one touching both and the first;
and so on around the ring. Sometimes the chain closes up exactly. If it does, would it still close up if you started with a different first
circle?
\end{think}

\section{Inversion in a circle}

\begin{definition}\label{def:inversion}
Let $\omega$ be a circle with centre $O$ and radius $r$. The \emph{inversion} in $\omega$ maps a point $P\ne O$ to the point $P'$ on the ray
$OP$ with $OP\cdot OP'=r^2$. In complex coordinates with $O=0$ it is $z\mapsto r^2/\bar z$.
\end{definition}

The inversion is an involution ($P''=P$), fixes $\omega$ pointwise and exchanges its inside and outside; points near $O$ go far away. It is
convenient to add one point $\infty$ to the plane and to say that $O$ and $\infty$ are exchanged. The plane with this point is the
\emph{extended plane} $\hat{\C}=\C\cup\{\infty\}$.

\begin{proposition}[distances]\label{prop:invdist}
For $P,Q\ne O$, \quad $P'Q'=\dfrac{r^2\,PQ}{OP\cdot OQ}$.
\end{proposition}
\begin{proof}
$\bigl|\frac{r^2}{\bar z}-\frac{r^2}{\bar w}\bigr|=r^2\frac{|\bar w-\bar z|}{|z|\,|w|}$.
\end{proof}

A \emph{generalised circle} is a circle or a line. Both are zero sets of equations
\begin{equation}\label{eq:gencircle}
\begin{gathered}
A|z|^2+Bz+\bar B\bar z+C=0,\\
A,C\in\R,\quad B\in\C,\quad |B|^2>AC .
\end{gathered}
\end{equation}
For $A\ne0$ this is the circle of centre $-\bar B/A$ and radius $\sqrt{|B|^2-AC}/|A|$ (complete the square:
$|z+\bar B/A|^2=(|B|^2-AC)/A^2$); for $A=0$ it is a line. The generalised circle passes through $0$ exactly when $C=0$.

\begin{theorem}[inversion and circles]\label{thm:invcircles}
Inversion maps generalised circles to generalised circles. More precisely, for inversion centred at $O$: lines through $O$ go to themselves;
lines not through $O$ and circles through $O$ are exchanged; circles not through $O$ go to circles not through $O$. Inversion preserves the
angle between two curves at a point where they meet (and reverses its orientation).
\end{theorem}
\begin{proof}
Scaling, take $r=1$ and $O=0$. If $w=1/\bar z$, then $z=1/\bar w$ and $|z|^2=1/|w|^2$. Substituting in \eqref{eq:gencircle} and multiplying by
$|w|^2$ gives
\[
A+B\,\bar w^{-1}|w|^2+\bar B\,w^{-1}|w|^2+C|w|^2=C|w|^2+Bw+\bar B\bar w+A=0,
\]
an equation of the same form with $A$ and $C$ exchanged. Reading off when $A$ or $C$ vanishes gives the four cases.

For angles: near a point $z\ne0$, $w=1/\bar z$ changes by $dw=-d\bar z/\bar z^2$ to first order, so a tangent vector $v$ at $z$ becomes
$-\bar v/\bar z^2$: it is conjugated, then multiplied by a fixed complex number. Conjugation reverses angles between vectors and
multiplication by a complex number preserves them, so the angle between two curves keeps its size and changes its sign.
\end{proof}

Since inversion is a bijection of $\hat{\C}$ that preserves generalised circles, it preserves tangency (two generalised circles are tangent
when they meet in exactly one point of $\hat{\C}$; parallel lines meet only at $\infty$). Two circles are \emph{orthogonal} if they cross at
right angles; by Pythagoras, circles of radii $\rho_1,\rho_2$ with centres $d$ apart are orthogonal exactly when $d^2=\rho_1^2+\rho_2^2$.

\begin{proposition}\label{prop:orthfixed}
A circle $\gamma$ orthogonal to $\omega$ is mapped to itself by the inversion in $\omega$. Conversely, a circle through a point $P$ and its
inverse $P'\ne P$ is orthogonal to $\omega$.
\end{proposition}
\begin{proof}
If $\gamma$ is orthogonal to $\omega$, the tangent from $O$ to $\gamma$ is a radius of $\omega$ and has length $r$, so the power of $O$ with
respect to $\gamma$ is $r^2$: every line through $O$ meeting $\gamma$ at $P$ and $Q$ has $OP\cdot OQ=r^2$, and $Q=P'$. (Here $O$ is outside
$\gamma$, so $P$ and $Q$ are on the same side of $O$.) Conversely, if $\gamma$ passes through $P$ and $P'$, the power of $O$ is $OP\cdot OP'=r^2$,
so the tangent from $O$ to $\gamma$ has length $r$ and ends on $\omega$, where it is perpendicular to the radius of $\gamma$: the circles are
orthogonal there.
\end{proof}

\section{Ptolemy's theorem}

Claudius Ptolemy used the following theorem around 150 to compute his tables of chords, the ancestors of our sine tables.

\begin{theorem}[Ptolemy's inequality]\label{thm:ptolemy}
For any four points $A,B,C,D$ of the plane,
\[
AB\cdot CD+AD\cdot BC\ \ge\ AC\cdot BD,
\]
with equality if and only if $A,B,C,D$ lie on a circle or a line, in this cyclic order.
\end{theorem}
\begin{proof}
Invert in a circle of radius $1$ centred at $A$. By the triangle inequality $B'C'+C'D'\ge B'D'$, and by Proposition~\ref{prop:invdist}
\[
\frac{BC}{AB\cdot AC}+\frac{CD}{AC\cdot AD}\ \ge\ \frac{BD}{AB\cdot AD}.
\]
Multiplying by $AB\cdot AC\cdot AD$ gives the inequality. Equality holds exactly when $C'$ lies on the segment $B'D'$, that is, when $B',C',D'$
are on a line (not through $A$, or through $A$) in this order; by Theorem~\ref{thm:invcircles} this means that $B,C,D$ lie on a circle through
$A$ (or on a line through $A$) with $C$ on the arc $BD$ not containing $A$.
\end{proof}

For a cyclic quadrilateral with sides $a,b,c,d$ and diagonals $e,f$, Ptolemy's theorem says $ac+bd=ef$. Applied to a rectangle it is
Pythagoras' theorem; applied to well-chosen quadrilaterals inscribed in a circle of diameter $1$, where chords are sines, it gives the
addition formula $\sin(\alpha+\beta)=\sin\alpha\cos\beta+\cos\alpha\sin\beta$, which is how Ptolemy used it.

\section{Steiner's porism}

Inversion does not preserve distances or radii, but it preserves one number attached to a pair of circles. For circles with radii
$\rho_1,\rho_2$ and centres $d$ apart, define the \emph{inversive distance}
\[
\delta=\frac{\rho_1^2+\rho_2^2-d^2}{2\rho_1\rho_2}.
\]
It is $1$ for internally tangent circles, $-1$ for externally tangent ones, $0$ for orthogonal ones, and $\delta=\cos\theta$ when the circles
cross at an angle $\theta$ (cosine law in the triangle formed by the two centres and an intersection point).

\begin{proposition}[inversive distance is invariant]\label{prop:invdistance}
If an inversion maps two circles (not through its centre) to two circles, it preserves $|\delta|$.
\end{proposition}
\begin{proof}
Attach to the circle with centre $c$ and radius $\rho$ the vector $X=(A,B,C)=\frac1\rho\bigl(1,-\bar c,|c|^2-\rho^2\bigr)$, so that the circle
is \eqref{eq:gencircle} divided by $\rho$. On such triples define the symmetric bilinear form
\[
\langle X_1,X_2\rangle=\operatorname{Re}\bigl(B_1\bar B_2\bigr)-\tfrac12\bigl(A_1C_2+A_2C_1\bigr).
\]
A direct computation gives $\langle X,X\rangle=\frac1{\rho^2}\bigl(|c|^2-(|c|^2-\rho^2)\bigr)=1$, and for two circles
\[
\langle X_1,X_2\rangle=\frac1{\rho_1\rho_2}\Bigl(\operatorname{Re}(c_1\bar c_2)-\tfrac12\bigl(|c_2|^2-\rho_2^2+|c_1|^2-\rho_1^2\bigr)\Bigr)=\delta .
\]
Similarities change $\rho_1$, $\rho_2$ and $d$ by the same factor, so they preserve $\delta$; and every inversion is a similarity followed by the
inversion $z\mapsto1/\bar z$ followed by a similarity. By the proof of Theorem~\ref{thm:invcircles}, $z\mapsto1/\bar z$ maps the circle with
triple $(A,B,C)$ to the one with triple $(C,B,A)$, up to a non-zero real factor (we must divide by the new radius). The form is unchanged by
exchanging $A$ and $C$ in both arguments, and a real factor $\lambda$ multiplies $\langle X,X\rangle=1$ by $\lambda^2$, so after renormalising
the factor is $\pm1$, and $\delta$ is preserved up to sign.
\end{proof}

The sign records orientation: it changes when an inversion turns a circle ``inside out'', as happens to a circle containing the centre of
inversion. For the problems below only $|\delta|$ matters.

\begin{lemma}[concentric position]\label{lem:concentric}
Two circles that do not meet can be mapped by an inversion to two concentric circles.
\end{lemma}
\begin{proof}
Put the centres at $0$ and $d>0$ on the real axis, with radii $\rho_1,\rho_2$. We look for two real points $p\ne q$ that are inverse with respect
to both circles: $pq=\rho_1^2$ and $(p-d)(q-d)=\rho_2^2$. Subtracting, $p+q=s$ with $s=(\rho_1^2+d^2-\rho_2^2)/d$, so $p,q$ are the roots of
$x^2-sx+\rho_1^2$, which are real and distinct when $s^2>4\rho_1^2$. If the first circle is inside the second, $d<\rho_2-\rho_1$, then
$\rho_2^2-\rho_1^2-d^2>2\rho_1d$ (equivalent to $\rho_2^2>(\rho_1+d)^2$), so $-s>2\rho_1$. If the circles are outside each other,
$d>\rho_1+\rho_2$, then $s-2\rho_1=\bigl((d-\rho_1)^2-\rho_2^2\bigr)/d>0$. Either way $p$ and $q$ exist.

Now invert in any circle centred at $p$. By Proposition~\ref{prop:orthfixed}, every circle through $p$ and $q$ is orthogonal to both given
circles; these circles go to lines through $q'$, the image of $q$. The images of the two given circles are orthogonal to all lines through
$q'$, so they are circles centred at $q'$.
\end{proof}

\begin{theorem}[Steiner's porism]\label{thm:steiner}
Let $\Gamma_1$ be a circle inside a circle $\Gamma_2$, not touching it. Start a chain of circles, each tangent to $\Gamma_1$ and $\Gamma_2$ and to the
previous one. If the chain closes up (the $n$-th circle touches the first) for one choice of the first circle, it closes up, with the same
number $n$ of circles and the same number of turns around $\Gamma_1$, for every choice.
\end{theorem}
\begin{proof}
By Lemma~\ref{lem:concentric}, an inversion maps $\Gamma_1$ and $\Gamma_2$ to concentric circles of radii $a<b$, and since it preserves tangency, it
maps chains to chains. Between concentric circles, every circle tangent to both has radius $\frac{b-a}2$ and centre on the circle of radius
$\frac{a+b}2$, and two consecutive ones are related by a rotation about the common centre by the angle $2\arcsin\frac{b-a}{b+a}$. The chain
closes after $n$ circles and $k$ turns exactly when $n$ times this angle is $2\pi k$, a condition that does not involve the first circle. Rotating
a closed chain gives a closed chain from any starting circle.
\end{proof}

\begin{figure}[ht]
\centering
\fig{ch03-steiner}
\caption{A Steiner chain of seven circles. It is the image, under an inversion, of a ring of seven equal circles between two concentric
circles, which is why it closes up wherever it starts.}
\label{fig:steiner}
\end{figure}

The word \emph{porism} comes from Euclid's lost book of \emph{Porisms}; it has come to mean a statement of the form ``if one solution exists, then
infinitely many do''. Poncelet's porism in Chapter~5 is another one.

\section{The Pappus chain}

The \emph{arbelos} (``shoemaker's knife'') is the region inside a half-disc with diameter $AB$ and outside the half-discs with diameters $AC$
and $CB$, for a point $C$ of $AB$. Archimedes studied it in his \emph{Book of Lemmas}. Around 300 Pappus of Alexandria described an infinite chain
of circles inside it and proved the following theorem.

\begin{theorem}[Pappus]\label{thm:pappus}
Let $C_0$ be the circle on $CB$, and let $C_1,C_2,\dots$ be circles in the arbelos, each tangent to the circles on $AB$ and on $AC$ and to the
previous one. Then the centre of $C_n$ is at height $n\,d_n$ above $AB$, where $d_n$ is the diameter of $C_n$.
\end{theorem}
\begin{proof}
Invert in the circle centred at $A$ that is orthogonal to $C_n$ (its radius is the length of the tangent from $A$ to $C_n$). The circles on
$AB$ and $AC$ pass through $A$ and have their centres on $AB$, so they become two lines perpendicular to $AB$ (the line $AB$ goes to itself and
angles are preserved). The circles $C_0,\dots,C_n$, tangent to both, become circles of equal diameter between two parallel lines, each tangent to
the next: a vertical stack. The image of $C_0$ is symmetric about $AB$, because $C_0$ is, so it is centred on $AB$; and $C_n$ is its own
image by Proposition~\ref{prop:orthfixed}. So the centre of $C_n$ is $n$ diameters above $AB$.
\end{proof}

\begin{figure}[ht]
\centering
\fig{ch03-pappus}
\caption{The Pappus chain in an arbelos. The centre of $C_n$ is exactly $n$ diameters of $C_n$ above the line $AB$ (red), as the computation
from the inverted picture confirms.}
\label{fig:pappus}
\end{figure}

\section{Descartes' circle theorem}

The triples of Proposition~\ref{prop:invdistance} do more than measure inversive distance: they turn tangent circles into vectors with a
fixed inner product, and Descartes' theorem becomes linear algebra. To handle circles that contain the others, we orient them: a circle
gets a \emph{signed curvature} $k=1/\rho$ if the other circles of the configuration are outside it and $k=-1/\rho$ if they are inside it, and
its triple is $X=k\,(1,-\bar c,|c|^2-\rho^2)$. Then the first coordinate of $X$ is $k$ and $\langle X,X\rangle=1$, and for two tangent circles
in a configuration (externally tangent with positive curvatures, or one inside the other with opposite signs) the computation of
Proposition~\ref{prop:invdistance} gives $\langle X_1,X_2\rangle=-1$.

\begin{think}
Three coins of radii $1$, $2$ and $3$ touch each other. What is the radius of the circle that encloses them and touches all three? (Look at the
triangle formed by their centres.)
\end{think}

\begin{theorem}[Descartes 1643, Soddy 1936; Lagarias, Mallows and Wilks 2002]\label{thm:descartes}
Let four circles be pairwise tangent at six distinct points, with signed curvatures $k_1,\dots,k_4$ and centres $c_1,\dots,c_4$. Then
\[
\Bigl(\sum_ik_i\Bigr)^2=2\sum_ik_i^2,\qquad\Bigl(\sum_ik_ic_i\Bigr)^2=2\sum_i(k_ic_i)^2,\qquad
\Bigl(\sum_ik_i\Bigr)\Bigl(\sum_ik_ic_i\Bigr)=2\sum_ik_i\,(k_ic_i).
\]
\end{theorem}
\begin{proof}
Write the triples in the real coordinates $A$, $\operatorname{Re}B$, $\operatorname{Im}B$, $C$ of $\R^4$; the form becomes $\langle X,Y\rangle=X^{\mathsf T}\Omega Y$ with
\[
\Omega=\begin{pmatrix}0&0&0&-\frac12\\0&1&0&0\\0&0&1&0\\-\frac12&0&0&0\end{pmatrix},\qquad
\Omega^{-1}=\begin{pmatrix}0&0&0&-2\\0&1&0&0\\0&0&1&0\\-2&0&0&0\end{pmatrix}.
\]
Let $M$ be the $4\times4$ matrix whose rows are the triples $X_1,\dots,X_4$. The Gram matrix $M\Omega M^{\mathsf T}$ has $1$ on the diagonal and $-1$
elsewhere: it is $G=2I-J$, where $J$ is the all-ones matrix. Since $J^2=4J$, one checks $G\bigl(\frac12I-\frac14J\bigr)=I$, so $G$ is invertible,
hence so is $M$, and
\[
\Omega^{-1}=M^{\mathsf T}G^{-1}M .
\]
For any two vectors $u,v$ (real or complex), $u^{\mathsf T}\Omega^{-1}v=(Mu)^{\mathsf T}G^{-1}(Mv)$. Take $u=v=e_1$, the first basis vector: $Me_1$
is the vector $k$ of curvatures, and $e_1^{\mathsf T}\Omega^{-1}e_1=0$, so
\[
0=k^{\mathsf T}\bigl(\tfrac12I-\tfrac14J\bigr)k=\tfrac12\sum k_i^2-\tfrac14\Bigl(\sum k_i\Bigr)^2 .
\]
This is the first formula. For the second, take $u=v=\varphi=(0,-1,i,0)$; then $M\varphi$ has entries $-\operatorname{Re}B_j+i\operatorname{Im}B_j=-\bar B_j=k_jc_j$,
and $\varphi^{\mathsf T}\Omega^{-1}\varphi=(-1)^2+i^2=0$. For the third, $e_1^{\mathsf T}\Omega^{-1}\varphi=0$.
\end{proof}

Given three mutually tangent circles, the first formula is a quadratic equation for $k_4$, whose two roots are the curvatures of the two
circles tangent to all three. The vectors give a sharper statement.

\begin{proposition}[the swap rule]\label{prop:swap}
Let $X_1,X_2,X_3,X_4$ be the vectors of four circles as in Theorem~\ref{thm:descartes}, and $X_4'$ that of the other circle tangent to the first
three. Then
\[
X_4+X_4'=2(X_1+X_2+X_3).
\]
In particular
\begin{equation}\label{eq:descartesswap}
\begin{aligned}
k_4+k_4'&=2(k_1+k_2+k_3),\\
k_4c_4+k_4'c_4'&=2(k_1c_1+k_2c_2+k_3c_3).
\end{aligned}
\end{equation}
\end{proposition}
\begin{proof}
The vectors $X$ with $\langle X,X_i\rangle=-1$ for $i=1,2,3$ form a line $Y+\R N$, where $Y=X_1+X_2+X_3$ (indeed $\langle Y,X_i\rangle=1-1-1=-1$) and $N\neq0$
is orthogonal to $X_1,X_2,X_3$. The Gram matrix of $X_1,X_2,X_3$ is $2I-J$ of size $3$, with eigenvalues $2,2,-1$; since the whole form has
signature $(3,1)$ (it is $|B|^2-AC$, and $-AC$ has one positive and one negative direction), the orthogonal complement of $X_1,X_2,X_3$ is
positive: $\langle N,N\rangle>0$. On the line, $\langle Y+tN,Y+tN\rangle=\langle Y,Y\rangle+t^2\langle N,N\rangle$ because $\langle Y,N\rangle=0$, so the
vectors of norm $1$ on it are $Y\pm t_0N$, symmetric about $Y$. Both $X_4$ and $X_4'$ are among them, and they are different circles, so
$X_4+X_4'=2Y$. Reading the first coordinate and the coordinate $-\bar B=kc$ gives \eqref{eq:descartesswap}.
\end{proof} For the coins of the question, $k=1,\frac12,\frac13$ and
$k_4=\frac{11}6\pm2\sqrt{\frac12+\frac16+\frac13}=\frac{11}6\pm2$: the enclosing circle has curvature $-\frac16$, radius $6$, and the small
circle in the middle has radius $\frac6{23}$.

\section{Apollonian gaskets}

Start with four mutually tangent circles, one enclosing the other three. In each of the four curvilinear triangles between them, put the
circle tangent to its three sides; this creates twelve new triangles; fill them too, and repeat forever. The result is an \emph{Apollonian
gasket}: a fractal whose holes have total area equal to that of the outer disc, and whose leftover set has Hausdorff dimension about
$1.3057$ (a remark: the dimension was estimated by McMullen and others, and is not known in closed form).

\begin{corollary}[integral gaskets]\label{cor:integral}
If the four starting curvatures are integers, all the curvatures in the gasket are integers. If moreover the starting curvatures times
the centres are Gaussian integers, so are all the $k_ic_i$.
\end{corollary}
\begin{proof}
Every new circle replaces one circle $k_4$ of a Descartes quadruple $(k_1,k_2,k_3,k_4)$ by the other root $k_4'=2(k_1+k_2+k_3)-k_4$ of
\eqref{eq:descartesswap}, and the new quadruple is again a Descartes quadruple. By induction on the generation, all are integers. The same
argument works for $k_4'c_4'=2(k_1c_1+k_2c_2+k_3c_3)-k_4c_4$.
\end{proof}

\begin{figure}[ht]
\centering
\fig{ch03-gasket}
\caption{The Apollonian gasket with root quadruple $(-1,2,2,3)$: the outer circle has radius $1$, the two large circles radius $\frac12$, and every
circle has an integer curvature. The circles of curvature $6$ come from $2\cdot(-1+2+3)-2$.}
\label{fig:gasket}
\end{figure}

Integral gaskets are a playground of number theory. In 2003 Ron Graham, Jeffrey Lagarias, Colin Mallows, Allan Wilks and Catherine Yan
asked which integers occur as curvatures; Jean Bourgain and Alex Kontorovich proved in 2014 that almost all integers that are not ruled out by
congruences do occur, and in 2023 Summer Haag, Clyde Kertzer, James Rickards and Katherine Stange found that some integers are missed for
subtler, quadratic-reciprocity reasons. These results are quoted as remarks; the notebook explores them numerically.

\section{M\"obius transformations}

A \emph{M\"obius transformation} is a map of $\hat{\C}$ of the form
\[
f(z)=\frac{az+b}{cz+d},\qquad a,b,c,d\in\C,\ ad-bc\ne0,
\]
with $f(-d/c)=\infty$ and $f(\infty)=a/c$ (and $f(\infty)=\infty$ if $c=0$). Composition of M\"obius maps corresponds to the product of their
matrices $\bigl(\begin{smallmatrix}a&b\\c&d\end{smallmatrix}\bigr)$, so they form a group, the image of $GL_2(\C)$, in which a matrix and its non-zero
multiples give the same map.

\begin{theorem}[M\"obius maps]\label{thm:mobius}
\begin{enumerate}
\item Every M\"obius map is a composition of translations, maps $z\mapsto\lambda z$, and $z\mapsto1/z$. It maps generalised circles to
generalised circles and preserves angles and their orientation.
\item For distinct $z_1,z_2,z_3\in\hat{\C}$ and distinct $w_1,w_2,w_3\in\hat{\C}$ there is exactly one M\"obius map with $f(z_j)=w_j$.
\item The \emph{cross-ratio}
\[
[z_1,z_2;z_3,z_4]=\frac{(z_1-z_3)(z_2-z_4)}{(z_1-z_4)(z_2-z_3)}
\]
is preserved by M\"obius maps, and it is real if and only if the four points lie on a generalised circle.
\end{enumerate}
\end{theorem}
\begin{proof}
(1) If $c=0$, $f$ is $z\mapsto(a/d)z+b/d$. If $c\ne0$, $f(z)=\frac ac-\frac{ad-bc}{c(cz+d)}$, a composition of $z\mapsto cz+d$, $z\mapsto1/z$,
$z\mapsto-\frac{ad-bc}cz$ and a translation. Translations and $z\mapsto\lambda z$ are similarities; $z\mapsto1/z$ is the inversion $z\mapsto1/\bar z$
followed by the reflection $z\mapsto\bar z$. Each of these preserves generalised circles (Theorem~\ref{thm:invcircles}), and the two orientation
reversals cancel.

(2) The map $g(z)=\frac{(z-z_1)(z_2-z_3)}{(z-z_3)(z_2-z_1)}$ (with the obvious changes if one of the $z_j$ is $\infty$) sends $z_1,z_2,z_3$ to
$0,1,\infty$. If $h$ does the same for the $w_j$, then $h^{-1}g$ works. For uniqueness, a M\"obius map fixing $0,1,\infty$ has $c=0$ (it fixes
$\infty$), $b=0$ (it fixes $0$) and $a=d$ (it fixes $1$), so it is the identity.

(3) Each of the generating maps of (1) preserves the cross-ratio (for translations and scalings the factors cancel; for $z\mapsto1/z$,
$\frac1{z_1}-\frac1{z_3}=\frac{z_3-z_1}{z_1z_3}$ and the products $z_1z_3\,z_2z_4$ and $z_1z_4\,z_2z_3$ cancel). Now apply the map $f$ sending
$z_2,z_3,z_4$ to $1,0,\infty$: the cross-ratio becomes $[f(z_1),1;0,\infty]=f(z_1)$. The four points are on a generalised circle if and only if
their images are, that is, if and only if $f(z_1)$ is on the generalised circle through $1,0,\infty$, the real line.
\end{proof}

\begin{remark}
M\"obius maps with $ad-bc=1$ are classified by $\operatorname{tr}^2=(a+d)^2$: \emph{elliptic} if $\operatorname{tr}^2\in[0,4)$ (conjugate to a rotation
$z\mapsto e^{i\theta}z$), \emph{parabolic} if $\operatorname{tr}^2=4$ and $f\ne\mathrm{id}$ (conjugate to $z\mapsto z+1$), and \emph{loxodromic} otherwise
(conjugate to $z\mapsto\lambda z$ with $|\lambda|\ne1$). Exercise~\ref{ex:trace} proves it. In Chapter~10 the M\"obius maps preserving the upper
half-plane, the real matrices of $SL_2(\R)$, become the isometries of the hyperbolic plane.
\end{remark}

\section{A gem: Feuerbach's theorem}

Chapter~1 promised a proof that the nine-point circle touches the incircle and the three excircles. Feuerbach's original proof was a long
computation; the classical modern proof inverts in the circle centred at the midpoint of $BC$ through the points where the incircle and the
$A$-excircle touch $BC$, which fixes both circles and maps the nine-point circle to their fourth common tangent (Coxeter and Greitzer,
Section~5.6). Here is a proof in the coordinates of Chapter~1, which shows where the tangency comes from.

\begin{theorem}[Feuerbach, 1822]\label{thm:feuerbach}
The nine-point circle of a triangle is tangent internally to its incircle and externally to its three excircles.
\end{theorem}
\begin{proof}
Use the coordinates of Proposition~\ref{prop:incentre}: the circumcircle is the unit circle, $a=u^2$, $b=v^2$, $c=w^2$, and the incentre is
$i=-(uv+vw+wu)$. Write $s_1=u+v+w$, $s_2=uv+vw+wu$, $s_3=uvw$; since $\bar u=1/u$ and so on, $\bar s_1=s_2/s_3$. The nine-point centre is
$n=\frac12(a+b+c)=\frac12(s_1^2-2s_2)$, so
\[
n-i=\tfrac12s_1^2,\qquad NI=\tfrac12|s_1|^2 .
\]
The inradius $r$ is the distance from $i$ to the line $bc$. By Lemma~\ref{lem:chord} it is $\frac12|b+c-i-bc\,\bar i|$. Here
$\bar i=-\bar s_2=-s_1/s_3$, so $bc\,\bar i=-vw\,s_1/u$, and
\[
u\bigl(b+c-i-bc\,\bar i\bigr)=u v^2+uw^2+u\,s_2+vw\,s_1=(u+v)(v+w)(w+u)=s_1s_2-s_3,
\]
as one checks by expanding. Since $s_1s_2-s_3=s_3\bigl(s_1\bar s_1-1\bigr)=s_3(|s_1|^2-1)$ and $|u|=|s_3|=1$, we get $r=\frac12\bigl|1-|s_1|^2\bigr|$.
Also $|i|=|s_2|=|\bar s_1s_3|=|s_1|$, and the incentre is inside the circumcircle, so $|s_1|<1$. Therefore
\[
NI=\tfrac12|s_1|^2=\tfrac12-r=\frac R2-r ,
\]
which says exactly that the incircle (radius $r$) touches the nine-point circle (radius $\frac R2$) from inside. On the way we have re-proved
Euler's formula $OI^2=|s_1|^2=1-2r$ of Exercise~\ref{ex:eulerOI}.

For the excircle opposite $A$, replace $u$ by $-u$. All the identities above are polynomial identities in $u,v,w$, so they still hold. The
point $i_a=-(-uv+vw-wu)$ lies on the line through $u^2$ and $-vw$ (the internal bisector from $A$), and on the lines through $v^2$ and $+wu$ and
through $w^2$ and $+uv$; since $+wu$ is the midpoint of the arc $CA$ \emph{containing} $B$, which is diametrically opposite to the midpoint of
the other arc, the line through $v^2$ and $+wu$ is the external bisector at $B$, and likewise at $C$. So $i_a$ is the excentre, the distance
from $i_a$ to $BC$ is the exradius $r_a=\frac12\bigl||s_1'|^2-1\bigr|$ with $s_1'=-u+v+w$, and $|i_a|=|s_1'|>1$ because the excentre is outside
the circumcircle. Hence $NI_a=\frac12|s_1'|^2=\frac12+r_a=\frac R2+r_a$: the excircle touches the nine-point circle from outside.
\end{proof}

\begin{beauty}
One transformation, two lines of algebra, and the configurations of tangent circles fall into place: Steiner's chain is a ring of equal
pearls seen through a distorting lens, Pappus' chain is a stack of coins, and Descartes' theorem is the statement that four vectors of norm
one, at mutual inner product $-1$, live in a space with a quadratic form of signature $(3,1)$, the form of Einstein's space-time. The
curvatures in an Apollonian gasket are integers forever, and which integers appear is a research problem that was still producing surprises
in 2023.
\end{beauty}

\section{Exercises}

\begin{exercise}[$\star$]\label{ex:invline}
Find the image of the line $\operatorname{Re}z=1$ under the inversion in the unit circle, and the image of the circle $|z-2|=1$.
\end{exercise}

\begin{exercise}[$\star$]\label{ex:vanschooten}
Let $ABC$ be equilateral and $P$ a point of the arc $BC$ of its circumcircle not containing $A$. Show that $PA=PB+PC$.
\end{exercise}

\begin{exercise}[$\star\star$]\label{ex:twins}
In an arbelos with $AC=2r_1$ and $CB=2r_2$, draw the perpendicular to $AB$ at $C$. Show that the two circles inscribed in the arbelos on
either side of it, each tangent to the perpendicular, to the big semicircle and to one of the small ones, have the same radius
$\frac{r_1r_2}{r_1+r_2}$ (Archimedes' twin circles).
\end{exercise}

\begin{exercise}[$\star\star$]\label{ex:steinerformula}
Show that a Steiner chain of $n$ circles winding once around exists between circles of radii $\rho<R$ with centres $d$ apart if and only if
$(R-\rho)^2-4R\rho\tan^2\frac\pi n=d^2$.
\end{exercise}

\begin{exercise}[$\star$]\label{ex:soddy}
Three circles of radii $2$, $3$ and $10$ touch each other externally. Find the radii of the two circles tangent to all three.
\end{exercise}

\begin{exercise}[$\star\star$]\label{ex:pappusradius}
In the arbelos with $AB=1$ and $AC=t$, show that the circle $C_n$ of the Pappus chain has radius $r_n=\dfrac{t(1-t)}{2\bigl(t+n^2(1-t)^2\bigr)}$.
\end{exercise}

\begin{exercise}[$\star$]\label{ex:gasketroot}
In the gasket of Figure~\ref{fig:gasket}, list the curvatures of the circles of the first two generations, and show that no curvature in it
is divisible by $4$. (Work modulo $4$ in \eqref{eq:descartesswap}.)
\end{exercise}

\begin{exercise}[$\star$]\label{ex:mobius3}
Find the M\"obius map sending $1,i,-1$ to $0,1,\infty$, and the image of the unit disc under it.
\end{exercise}

\begin{exercise}[$\star\star$]\label{ex:trace}
Let $f$ be a M\"obius map with matrix of determinant $1$, $f\ne\mathrm{id}$. Show that $f$ has one or two fixed points in $\hat{\C}$, that it has one if
and only if $\operatorname{tr}^2=4$, in which case it is conjugate to $z\mapsto z+1$, and that otherwise it is conjugate to $z\mapsto\lambda z$ with
$\lambda+\lambda^{-1}+2=\operatorname{tr}^2$.
\end{exercise}

\begin{exercise}[$\star\star$]\label{ex:stereo}
Inversion in a sphere of $\R^3$ is defined as in the plane. Show that it maps spheres and planes to spheres and planes, and deduce that the
stereographic projection from the north pole $N$ of the unit sphere to the equatorial plane maps circles on the sphere to circles or lines.
(Stereographic projection is the restriction of the inversion in the sphere of centre $N$ and radius $\sqrt2$.)
\end{exercise}

\begin{exercise}[$\star$]\label{ex:orthcircles}
Let $P\ne Q$ be two points. Show that every circle through $P$ and $Q$ is orthogonal to every circle with respect to which $P$ and $Q$ are
inverse.
\end{exercise}

\begin{exercise}[$\star\star\star$]\label{ex:hexlet}
Soddy's hexlet: between three mutually tangent spheres, a chain of spheres each tangent to the three and to the previous one always closes
up after exactly six spheres. Prove it by inverting at the point of tangency of two of the three spheres.
\end{exercise}

\section{Hints}

\begin{hint}{ex:invline}
Use Theorem~\ref{thm:invcircles}; find the images of two or three points.
\end{hint}
\begin{hint}{ex:vanschooten}
Ptolemy for $ABPC$.
\end{hint}
\begin{hint}{ex:twins}
Coordinates: $C=0$, $A=-2r_1$, $B=2r_2$. Write the two tangency conditions for a circle of radius $x$ centred at $(-x,y)$.
\end{hint}
\begin{hint}{ex:steinerformula}
The inversive distance of the two circles is invariant; compute it for the concentric picture, where $\frac{b-a}{b+a}=\sin\frac\pi n$.
\end{hint}
\begin{hint}{ex:soddy}
Theorem~\ref{thm:descartes}.
\end{hint}
\begin{hint}{ex:pappusradius}
Invert at $A$ with radius $1$ and invert the circle of the stack back, using the image of a circle under $z\mapsto1/\bar z$.
\end{hint}
\begin{hint}{ex:gasketroot}
The root quadruple modulo $4$ is $(3,2,2,3)$; check that the set of quadruples reachable modulo $4$ never contains $0$.
\end{hint}
\begin{hint}{ex:mobius3}
Theorem~\ref{thm:mobius}(2).
\end{hint}
\begin{hint}{ex:trace}
Fixed points solve $cz^2+(d-a)z-b=0$. Conjugate to send the fixed points to $0$ and $\infty$, or the single one to $\infty$.
\end{hint}
\begin{hint}{ex:stereo}
In $\R^3$, a sphere or plane is $A|x|^2+\langle B,x\rangle+C=0$.
\end{hint}
\begin{hint}{ex:orthcircles}
Proposition~\ref{prop:orthfixed}.
\end{hint}
\begin{hint}{ex:hexlet}
After the inversion, two of the spheres become parallel planes and the third a sphere between them; the chain becomes a ring of equal spheres
around it.
\end{hint}

\section{Solutions}

\begin{solution}{ex:invline}
\idea A line not through the centre becomes a circle through the centre; find its diameter from the closest point.
\writeup The point $1$ is fixed and the point at infinity goes to $0$, so the image is a circle through $0$ and $1$, symmetric about the real axis:
the circle with diameter $[0,1]$, that is $|z-\frac12|=\frac12$. The circle $|z-2|=1$ is symmetric about the real axis and meets it at $1$ and $3$,
which go to $1$ and $\frac13$: the image is the circle with diameter $[\frac13,1]$, centre $\frac23$, radius $\frac13$.
\end{solution}

\begin{solution}{ex:vanschooten}
\idea A cyclic quadrilateral with three equal sides is made for Ptolemy.
\writeup $A,B,P,C$ lie on the circle in this order, so Ptolemy's equality gives $AB\cdot PC+AC\cdot BP=AP\cdot BC$. Dividing by the common side
length, $PC+PB=PA$.
\end{solution}

\begin{solution}{ex:twins}
\idea Each twin is tangent to two circles and a line, and two tangency conditions give two equations for the centre.
\writeup With $C=0$, the big circle has centre $r_2-r_1$ and radius $r_1+r_2$, the left small circle centre $-r_1$ and radius $r_1$. The left twin,
of radius $x$, is tangent to the line $\operatorname{Re}z=0$, so its centre is $(-x,y)$. Tangency to the small circle (externally) and to the big
circle (internally) gives $(r_1-x)^2+y^2=(r_1+x)^2$ and $(r_2-r_1+x)^2+y^2=(r_1+r_2-x)^2$. The first gives $y^2=4r_1x$. Substituting in the second,
$(r_2-r_1)^2+2(r_2-r_1)x+x^2+4r_1x=(r_1+r_2)^2-2(r_1+r_2)x+x^2$, so $(4r_2+4r_1)x=4r_1r_2$ and $x=\frac{r_1r_2}{r_1+r_2}$. The formula is symmetric
in $r_1,r_2$, so the right twin has the same radius.
\end{solution}

\begin{solution}{ex:steinerformula}
\idea The closing condition is simple for concentric circles, and inversive distance carries it to any pair.
\writeup By Lemma~\ref{lem:concentric} and Proposition~\ref{prop:invdistance}, the given circles have the same $|\delta|$ as concentric circles of radii
$a<b$, namely $\frac{a^2+b^2}{2ab}$, and by the proof of Theorem~\ref{thm:steiner} a chain of $n$ circles winding once exists if and only if
$\frac{b-a}{b+a}=\sin\frac\pi n=:s$. Then $\frac ba=\frac{1+s}{1-s}$ and $\frac{a^2+b^2}{2ab}=\frac{(1-s)^2+(1+s)^2}{2(1-s^2)}=\frac{1+s^2}{1-s^2}=1+2\tan^2\frac\pi n$.
(As $\frac{a^2+b^2}{2ab}$ increases with $\frac ba>1$, this value determines $\frac ba$.) For the given circles, one inside the other, $\delta=
\frac{R^2+\rho^2-d^2}{2R\rho}>1$, so the condition is $R^2+\rho^2-d^2=2R\rho+4R\rho\tan^2\frac\pi n$, that is $(R-\rho)^2-4R\rho\tan^2\frac\pi n=d^2$.
\end{solution}

\begin{solution}{ex:soddy}
\idea Descartes' quadratic equation, with the two roots for the two solutions.
\writeup $k_1+k_2+k_3=\frac12+\frac13+\frac1{10}=\frac{14}{15}$ and $k_1k_2+k_2k_3+k_3k_1=\frac16+\frac1{30}+\frac1{20}=\frac14$. Solving the first formula of
Theorem~\ref{thm:descartes} for $k_4$ gives $k_4=k_1+k_2+k_3\pm2\sqrt{k_1k_2+k_2k_3+k_3k_1}=\frac{14}{15}\pm1$: $k_4=\frac{29}{15}$ (radius
$\frac{15}{29}$, the small circle in the middle) or $k_4=-\frac1{15}$ (radius $15$, the circle enclosing all three).
\end{solution}

\begin{solution}{ex:pappusradius}
\idea In the inverted picture the stack is explicit; bring one circle back.
\writeup Inversion in the unit circle at $A=0$ maps the circles on $AB$ and $AC$ (through $0$, centres $\frac12$ and $\frac t2$) to the lines
$\operatorname{Re}z=1$ and $\operatorname{Re}z=\frac1t$, and the circle on $CB$ (through $t$ and $1$) to the circle on $[1,\frac1t]$. The image of $C_n$ is the
circle of radius $\frac w2$, $w=\frac1t-1$, centred at $p=\frac12(1+\frac1t)+inw$. A circle of centre $p$ and radius $s$ not through $0$ goes under
$z\mapsto1/\bar z$ to the circle of centre $p/(|p|^2-s^2)$ and radius $s/\bigl||p|^2-s^2\bigr|$ (check the two points of the circle on the line $0p$).
Here $|p|^2-\frac{w^2}4=\frac14\bigl((1+\frac1t)^2-(\frac1t-1)^2\bigr)+n^2w^2=\frac1t+\frac{n^2(1-t)^2}{t^2}$, and
\[
r_n=\frac{(1-t)/(2t)}{\frac1t+\frac{n^2(1-t)^2}{t^2}}=\frac{t(1-t)}{2\bigl(t+n^2(1-t)^2\bigr)} .
\]
The height of the centre is $\operatorname{Im}\frac p{|p|^2-s^2}=\frac{nw}{|p|^2-s^2}=2nr_n$, confirming Theorem~\ref{thm:pappus}.
\end{solution}

\begin{solution}{ex:gasketroot}
\idea New curvatures are $2(k_1+k_2+k_3)-k_4$, a linear rule that can be followed modulo $4$.
\writeup From $(-1,2,2,3)$: replacing $-1$ gives $2(7)+1=15$; replacing a $2$ gives $2(4)-2=6$; replacing $3$ gives $2(3)-3=3$. So the first
generation has curvatures $15,6,6,3$ (the second $3$ is the bottom circle of the figure). The second generation, from the quadruples
$(15,2,2,3)$, $(-1,6,2,3)$, $(-1,2,6,3)$, $(-1,2,2,3)$ by swapping each entry except the one just created, has the
curvatures $6$, $11$, $14$, $15$, $23$, $35$ and $38$ (several of them more than once, for symmetric circles). Modulo $4$ the root is $(3,2,2,3)$. The
rule $k\mapsto2(k_1+k_2+k_3)-k$ modulo $4$ only depends on $2(k_1+k_2+k_3)\bmod4$, that is on the parity of $k_1+k_2+k_3$. In the root, every
quadruple modulo $4$ has two entries $\equiv2$ and two odd entries. Replacing an odd entry $k$: the other three are $2,2$ and an odd number, with odd
sum, so $2(k_1+k_2+k_3)\equiv2$ and the new entry is $\equiv2-k$, odd. Replacing an entry $2$: the other three sum to $2+$odd$+$odd, which is
even, so the new entry is $\equiv-2\equiv2$. So every
quadruple has two entries $\equiv2\pmod 4$ and two odd entries, and no curvature is divisible by $4$.
\end{solution}

\begin{solution}{ex:mobius3}
\idea The proof of Theorem~\ref{thm:mobius}(2) gives the formula; the unit circle goes to the generalised circle through the images, the real line.
\writeup $f(z)=\frac{(z-1)(i+1)}{(z+1)(i-1)}=-i\,\frac{z-1}{z+1}$. Check: $f(1)=0$, $f(-1)=\infty$, $f(i)=-i\frac{i-1}{i+1}=-i\cdot i=1$. The unit circle goes
to the real line, and $f(0)=i$, so the unit disc goes to the upper half-plane.
\end{solution}

\begin{solution}{ex:trace}
\idea Move the fixed points to standard places; conjugation does not change the trace.
\writeup A fixed point $z\ne\infty$ satisfies $cz^2+(d-a)z-b=0$, and $\infty$ is fixed exactly when $c=0$. If $c\ne0$, this quadratic has
discriminant $(d-a)^2+4bc=(a+d)^2-4$, so two fixed points if $\operatorname{tr}^2\ne4$ and one if $\operatorname{tr}^2=4$. If $c=0$, $f(z)=\frac adz+\frac bd$ with
$ad=1$: fixed points $\infty$ and, if $a\ne d$, one more; $a=d$ means $a=d=\pm1$, $\operatorname{tr}^2=4$, and then $b\ne0$. The trace (up to sign) is
invariant under conjugation by M\"obius maps of determinant $1$. With one fixed point, conjugate it to $\infty$: $f(z)=z+\beta$, $\beta\ne0$, and
conjugating by $z\mapsto z/\beta$ gives $z+1$. With two, send them to $0$ and $\infty$: $f(z)=\lambda z$, with normalised matrix
$\operatorname{diag}(\lambda^{1/2},\lambda^{-1/2})$, so $\operatorname{tr}^2=\lambda+\lambda^{-1}+2$.
\end{solution}

\begin{solution}{ex:stereo}
\idea Repeat the computation of Theorem~\ref{thm:invcircles} in $\R^3$, then cut by the unit sphere.
\writeup With centre $0$ and radius $1$, $y=x/|x|^2$, $x=y/|y|^2$, $|x|^2=1/|y|^2$, and $A|x|^2+\langle B,x\rangle+C=0$ becomes, after multiplying by $|y|^2$,
$A+\langle B,y\rangle+C|y|^2=0$: spheres and planes go to spheres and planes. Now let $N=(0,0,1)$ and let $\iota$ be the inversion centred at $N$ with
radius $\sqrt2$: $\iota(x)=N+\frac{2(x-N)}{|x-N|^2}$, a point of the line $Nx$. A point $x$ of the unit sphere satisfies $|x-N|^2=2-2x_3$, so the
third coordinate of $\iota(x)$ is $1+\frac{2(x_3-1)}{2-2x_3}=0$: $\iota$ maps the sphere into the plane $x_3=0$, along the lines through $N$, which is
the definition of the stereographic projection. A
circle on the sphere is its intersection with a plane, and its image is the intersection of the plane $x_3=0$ with the image of that plane, a sphere
or a plane: a circle or a line.
\end{solution}

\begin{solution}{ex:orthcircles}
\idea This is the converse half of Proposition~\ref{prop:orthfixed}, applied to each circle in turn.
\writeup Let $\omega$ be a circle with respect to which $P$ and $Q=P'$ are inverse. A circle through $P$ and $P'$ is orthogonal to $\omega$ by
Proposition~\ref{prop:orthfixed}. (The line $PQ$ passes through the centre of $\omega$ and is orthogonal to it too.)
\end{solution}

\begin{solution}{ex:hexlet}
\idea Inversion turns two tangent spheres into parallel planes, and then everything is a problem about equal balls.
\writeup Let $S_1,S_2,S_3$ be the spheres, and invert at the point where $S_1$ and $S_2$ touch. They become two parallel planes $\Pi_1,\Pi_2$, at
distance $h$ say, and $S_3$, tangent to both, becomes a ball of diameter $h$ between them. Every sphere of the chain is tangent to $\Pi_1$, $\Pi_2$
and the ball $S_3'$, so it is a ball of diameter $h$ between the planes whose centre is at distance $h$ from the centre of $S_3'$, in the middle
plane. Consecutive balls of the chain are tangent, so their centres are at distance $h$ from each other. In the middle plane, the centres lie on
the circle of radius $h$ around the centre of $S_3'$, consecutive ones at distance $h$: they form a regular hexagon, so the chain closes after
exactly six balls, wherever it starts. Inverting back gives Soddy's hexlet.
\end{solution}

\section*{Sources and further reading}

\begin{itemize}
\item R.~Descartes, letter to Princess Elisabeth of November 1643, in \emph{\OE uvres de Descartes}, ed.\ Adam and Tannery, vol.~4; F.~Soddy,
\emph{The Kiss Precise}, Nature 137 (1936), 1021 (the epigraph is its second stanza).
\item J.~C.~Lagarias, C.~L.~Mallows and A.~R.~Wilks, \emph{Beyond the Descartes circle theorem}, American Mathematical Monthly 109 (2002), 338--361
(the vector proof and the complex Descartes theorem).
\item R.~L.~Graham, J.~C.~Lagarias, C.~L.~Mallows, A.~R.~Wilks and C.~H.~Yan, \emph{Apollonian circle packings: number theory}, Journal of Number
Theory 100 (2003), 1--45; J.~Bourgain and A.~Kontorovich, \emph{On the local-global conjecture for integral Apollonian gaskets}, Inventiones
Mathematicae 196 (2014), 589--650; S.~Haag, C.~Kertzer, J.~Rickards and K.~E.~Stange, \emph{The local-global conjecture for Apollonian circle
packings is false}, Annals of Mathematics 200 (2024).
\item C.~T.~McMullen, \emph{Hausdorff dimension and conformal dynamics III: computation of dimension}, American Journal of Mathematics 120 (1998),
691--721 (the dimension $1.3057$ of the gasket).
\item D.~Pedoe, \emph{Geometry: A Comprehensive Course}, Dover, 1988 (inversion, inversive distance, Steiner's porism); H.~S.~M.~Coxeter, \emph{The
inversive plane and hyperbolic space}, Abhandlungen aus dem Mathematischen Seminar der Universit\"at Hamburg 29 (1966), 217--242.
\item H.~S.~M.~Coxeter and S.~L.~Greitzer, \emph{Geometry Revisited}, MAA, 1967, Chapter~5 (inversion, Feuerbach by inversion, Steiner's porism).
\item T.~Needham, \emph{Visual Complex Analysis}, Oxford University Press, 1997, Chapter~3 (M\"obius maps and inversion, beautifully drawn).
\item L.~Bankoff, \emph{Are the twin circles of Archimedes really twins?}, Mathematics Magazine 47 (1974), 214--218; Pappus, \emph{Collection},
Book~IV (the chain in the arbelos).
\item W.~Thomson, \emph{Extrait d'une lettre de M.~William Thomson \`a M.~Liouville}, Journal de Math\'ematiques Pures et Appliqu\'ees 10 (1845),
364--367 (inversion in electrostatics, the method of images).
\end{itemize}
